% Figure: forest, leaf blocks and levels (Chapter 2).
% Geometry and tree topology: generated/forest.tex, written by generate.jl
% from its BlockTreeMesh fixture; labels and annotations are hand-written.
\begin{figure}
  \centering
  \begin{subfigure}{\textwidth}
    \centering
    \begin{tikzpicture}[x=0.98cm,y=0.98cm]
      \def\fxnode#1#2#3#4#5#6{}%
      \def\fxedge#1#2{}%
      \def\fxleaf#1#2#3#4#5#6#7#8{%
        \path[tr level #2] (#5,#6) rectangle (#7,#8);
        \pgfmathtruncatemacro\fxnx{\fxBlockCellsX-1}%
        \pgfmathtruncatemacro\fxny{\fxBlockCellsY-1}%
        \foreach \q in {1,...,\fxnx}
          \draw[tr cell] ({#5+\q*(#7-#5)/\fxBlockCellsX},#6) -- ++(0,#8-#6);
        \foreach \q in {1,...,\fxny}
          \draw[tr cell] (#5,{#6+\q*(#8-#6)/\fxBlockCellsY}) -- ++(#7-#5,0);
        \draw[tr block] (#5,#6) rectangle (#7,#8);
        \node[tr small,fill=white,rounded corners=0.6pt]
          at ({(#5+#7)/2},{(#6+#8)/2}) {#2,#3,#4};
      }%
      % Generated by docs/techref/figures/generate.jl; do not edit.
\def\fxBlockCellsX{2}\def\fxBlockCellsY{2}
\def\fxExtentX{12}\def\fxExtentY{4}
\fxleaf{1}{0}{0}{0}{0}{0}{4}{4}
\fxleaf{2}{1}{2}{0}{4}{0}{6}{2}
\fxleaf{3}{1}{3}{0}{6}{0}{8}{2}
\fxleaf{4}{1}{4}{0}{8}{0}{10}{2}
\fxleaf{5}{1}{5}{0}{10}{0}{12}{2}
\fxleaf{6}{1}{2}{1}{4}{2}{6}{4}
\fxleaf{7}{1}{3}{1}{6}{2}{8}{4}
\fxleaf{8}{1}{5}{1}{10}{2}{12}{4}
\fxleaf{9}{2}{8}{2}{8}{2}{9}{3}
\fxleaf{10}{2}{9}{2}{9}{2}{10}{3}
\fxleaf{11}{2}{8}{3}{8}{3}{9}{4}
\fxleaf{12}{2}{9}{3}{9}{3}{10}{4}
\fxnode{n-0-0-0}{0}{0}{0}{0}{1}
\fxnode{n-0-1-0}{3.1}{0}{1}{0}{0}
\fxnode{n-0-2-0}{8.825}{0}{2}{0}{0}
\fxnode{n-1-2-0}{1.6}{1}{2}{0}{1}
\fxnode{n-1-3-0}{2.6}{1}{3}{0}{1}
\fxnode{n-1-2-1}{3.6}{1}{2}{1}{1}
\fxnode{n-1-3-1}{4.6}{1}{3}{1}{1}
\fxnode{n-1-4-0}{6.2}{1}{4}{0}{1}
\fxnode{n-1-5-0}{7.2}{1}{5}{0}{1}
\fxnode{n-1-4-1}{9.7}{1}{4}{1}{0}
\fxnode{n-1-5-1}{12.2}{1}{5}{1}{1}
\fxnode{n-2-8-2}{8.2}{2}{8}{2}{1}
\fxnode{n-2-9-2}{9.2}{2}{9}{2}{1}
\fxnode{n-2-8-3}{10.2}{2}{8}{3}{1}
\fxnode{n-2-9-3}{11.2}{2}{9}{3}{1}
\fxedge{n-0-1-0}{n-1-2-0}
\fxedge{n-0-1-0}{n-1-3-0}
\fxedge{n-0-1-0}{n-1-2-1}
\fxedge{n-0-1-0}{n-1-3-1}
\fxedge{n-0-2-0}{n-1-4-0}
\fxedge{n-0-2-0}{n-1-5-0}
\fxedge{n-0-2-0}{n-1-4-1}
\fxedge{n-0-2-0}{n-1-5-1}
\fxedge{n-1-4-1}{n-2-8-2}
\fxedge{n-1-4-1}{n-2-9-2}
\fxedge{n-1-4-1}{n-2-8-3}
\fxedge{n-1-4-1}{n-2-9-3}

      % Roots (permanent level-zero blocks).
      \foreach \r in {0,1,2}
        \draw[tr root] ({4*\r},0) rectangle ({4*\r+4},4);
      \foreach \r/\note in {0/{},1/{refined by 2:1 balancing},2/{refined by \code{refine}}}
        \node[tr small,anchor=south,align=center] at ({4*\r+2},4.28) {root $(0,\r,0)$\\ \note};
      % Walls and periodic seam.
      \fill[tr wall] (0,-0.22) rectangle (12,0);
      \fill[tr wall] (0,4) rectangle (12,4.22);
      \draw (0,0) -- (12,0) (0,4) -- (12,4);
      \node[tr small,anchor=north] at (2,-0.25) {wall $\subset\physbnd$};
      \draw[tr thin vector,<->] (-0.35,1.2) -- (-0.35,2.8);
      \node[tr small,rotate=90,anchor=south] at (-0.38,2) {periodic};
      \draw[tr thin vector,<->] (12.35,1.2) -- (12.35,2.8);
      % Corner contact between levels 1 and 2.
      \node[tr vertex] (corner) at (8,2) {};
      \draw[black!70,line width=0.5pt] (corner) -- (8.5,-0.62)
        node[tr small,anchor=west] {corner contact, $\abs{\blocklevel{\blk}-\blocklevel{\blk'}}=1$};
      % Level legend.
      \begin{scope}[shift={(0.6,-1.3)}]
        \foreach \l/\w [count=\n from 0] in {0/{\basewidth{k}},1/{\basewidth{k}/2},2/{\basewidth{k}/4}} {
          \draw[tr block,tr level \l] ({3.9*\n},0) rectangle ++(0.3,0.3);
          \node[tr small,anchor=west] at ({3.9*\n+0.38},0.15)
            {$\lev=\l$: $\meshwidth{\lev}{k}=\w$};
        }
      \end{scope}
    \end{tikzpicture}
    \caption{Leaf blocks $\leaves$ in $\domain$, labelled with their keys
      $(\lev,i,j)$; thin lines are cell boundaries.}
    \label{fig:mesh:forest:leaves}
  \end{subfigure}

  \bigskip
  \begin{subfigure}{\textwidth}
    \centering
    \begin{tikzpicture}[x=0.9cm,y=-1.0cm]
      \def\fxleaf#1#2#3#4#5#6#7#8{}%
      \def\fxnode#1#2#3#4#5#6{%
        \ifnum#6=1
          \node[draw=black,tr level #3,minimum size=7pt,inner sep=0pt] (#1) at (#2,#3) {};
          \node[tr tiny,anchor=north] at (#1.south) {#3,#4,#5};
        \else
          \node[circle,draw=black,fill=white,minimum size=7pt,inner sep=0pt] (#1) at (#2,#3) {};
          \node[tr tiny,anchor=south] at ([yshift=1pt]#1.north) {#3,#4,#5};
        \fi
      }%
      \def\fxedge#1#2{\draw[black!70] (#1) -- (#2);}%
      % Generated by docs/techref/figures/generate.jl; do not edit.
\def\fxBlockCellsX{2}\def\fxBlockCellsY{2}
\def\fxExtentX{12}\def\fxExtentY{4}
\fxleaf{1}{0}{0}{0}{0}{0}{4}{4}
\fxleaf{2}{1}{2}{0}{4}{0}{6}{2}
\fxleaf{3}{1}{3}{0}{6}{0}{8}{2}
\fxleaf{4}{1}{4}{0}{8}{0}{10}{2}
\fxleaf{5}{1}{5}{0}{10}{0}{12}{2}
\fxleaf{6}{1}{2}{1}{4}{2}{6}{4}
\fxleaf{7}{1}{3}{1}{6}{2}{8}{4}
\fxleaf{8}{1}{5}{1}{10}{2}{12}{4}
\fxleaf{9}{2}{8}{2}{8}{2}{9}{3}
\fxleaf{10}{2}{9}{2}{9}{2}{10}{3}
\fxleaf{11}{2}{8}{3}{8}{3}{9}{4}
\fxleaf{12}{2}{9}{3}{9}{3}{10}{4}
\fxnode{n-0-0-0}{0}{0}{0}{0}{1}
\fxnode{n-0-1-0}{3.1}{0}{1}{0}{0}
\fxnode{n-0-2-0}{8.825}{0}{2}{0}{0}
\fxnode{n-1-2-0}{1.6}{1}{2}{0}{1}
\fxnode{n-1-3-0}{2.6}{1}{3}{0}{1}
\fxnode{n-1-2-1}{3.6}{1}{2}{1}{1}
\fxnode{n-1-3-1}{4.6}{1}{3}{1}{1}
\fxnode{n-1-4-0}{6.2}{1}{4}{0}{1}
\fxnode{n-1-5-0}{7.2}{1}{5}{0}{1}
\fxnode{n-1-4-1}{9.7}{1}{4}{1}{0}
\fxnode{n-1-5-1}{12.2}{1}{5}{1}{1}
\fxnode{n-2-8-2}{8.2}{2}{8}{2}{1}
\fxnode{n-2-9-2}{9.2}{2}{9}{2}{1}
\fxnode{n-2-8-3}{10.2}{2}{8}{3}{1}
\fxnode{n-2-9-3}{11.2}{2}{9}{3}{1}
\fxedge{n-0-1-0}{n-1-2-0}
\fxedge{n-0-1-0}{n-1-3-0}
\fxedge{n-0-1-0}{n-1-2-1}
\fxedge{n-0-1-0}{n-1-3-1}
\fxedge{n-0-2-0}{n-1-4-0}
\fxedge{n-0-2-0}{n-1-5-0}
\fxedge{n-0-2-0}{n-1-4-1}
\fxedge{n-0-2-0}{n-1-5-1}
\fxedge{n-1-4-1}{n-2-8-2}
\fxedge{n-1-4-1}{n-2-9-2}
\fxedge{n-1-4-1}{n-2-8-3}
\fxedge{n-1-4-1}{n-2-9-3}

      \node[tr small,anchor=east] at (-0.3,0) {$\lev=0$};
      \node[tr small,anchor=east] at (-0.3,1) {$\lev=1$};
      \node[tr small,anchor=east] at (-0.3,2) {$\lev=2$};
    \end{tikzpicture}
    \caption{The forest $\forest$: three trees, one per root. Squares are leaves,
      circles refined blocks.}
    \label{fig:mesh:forest:trees}
  \end{subfigure}
  \caption{Forest of the code fixture \code{BlockTreeMesh(base_mesh=(6,2),
    block_size=(2,2), extent=(12,4), periodic=(true,false), max_level=2)} in
    which \code{refine} marks only the keys $(0,2,0)$ and $(1,4,1)$. Root
    $(0,1,0)$ is refined by the 2:1 balance~\cref{eq:mesh:grading}, which also
    applies across corners: block $(1,3,0)$ touches the level-two block
    $(2,8,2)$ only at the marked vertex. Here $\basewidth{k}=2$, so the cell
    widths are $2$, $1$ and $1/2$. Drawn from
    \code{figures/generated/forest.tex}.}
  \label{fig:mesh:forest}
\end{figure}
